\documentclass[11pt]{article}
\usepackage[T1]{fontenc}
\usepackage{lmodern,amsmath,amssymb,amsthm,booktabs,microtype,xcolor}
\usepackage[margin=1in]{geometry}
\usepackage[colorlinks=true,linkcolor=teal,urlcolor=teal,citecolor=teal]{hyperref}
\newtheorem{theorem}{Theorem}
\newtheorem{lemma}[theorem]{Lemma}
\newtheorem{corollary}[theorem]{Corollary}
\newcommand{\F}{\mathbb F}
\newcommand{\m}{\mathfrak m}
\newcommand{\ord}{\operatorname{ord}}
\newcommand{\Fr}{\operatorname{Fr}}
\title{Frobenius rank and progression-free density\\over finite rings of characteristic two}
\author{Research note accompanying the ProveIt analysis}
\date{}
\begin{document}\maketitle
\begin{abstract}
For a finite local ring $A$ of characteristic two with $\m^3=0$, write $Q=|A|$, $q=|A/\m|$, and $s=|\{a^2:a\in\m\}|$. The loss of extremal upper density on forbidding polynomial geometric progressions is between $(1-1/q)^2s/Q^2$ and $s/Q^2$. For growing products of such rings the limiting extremal density is positive exactly when $\sum_i s_i/Q_i^2$ converges. The same criterion holds for unequal positive gaps in a fixed bounded range of powers. Rings with identical residue fields and identical dimensions in every radical layer can give opposite answers, even when their multiplication maps onto the second layer both have full image. This answers the characteristic-two, cube-zero case of a question in the ProveIt article and gives a family-level answer to Question 6.5 of Asada and coauthors.
\end{abstract}

\section{What changes when the radical has nonzero square?}
In a field, multiplication by a nonconstant polynomial increases degree. Over a finite local ring its highest coefficients can be nilpotent, and degree can decrease after cancellation. This matters for the question of how much of a large degree ball can be kept while avoiding three terms in geometric progression. Asada, Fourakis, Manski, McNew, Miller and Moreland asked how this problem changes on replacing the finite coefficient field by a finite ring \cite[Question 6.5]{asada}. The ProveIt article treated products of rings whose maximal ideals have square zero and asked for the invariant at nilpotence index three \cite{repo}. We treat all finite commutative local rings of characteristic two with cube-zero maximal ideal. The proof is unconditional.

Put $B_n(A)=\{f\in A[x]:\deg f\le n\}$, including zero. Give zero degree $-\infty$. For an integer $L\ge2$ call $E\subset A[x]$ admissible if it contains no three distinct terms
\[
              f,\quad r^a f,\quad r^{a+b}f
              \qquad(a,b\ge1,\ a+b\le L)
\]
with $r$ a nonunit of $A[x]$. Write $\Delta_L(A)$ for the supremum of $\limsup_n |E\cap B_n(A)|/|B_n(A)|$ over such $E$. Ordinary geometric progressions correspond to $L=2$. Use these definitions also for finite product rings, with the same degree bound in every component.

The roots of $X^q-X$ in $A$ form a coefficient field $k\simeq A/\m$. Indeed each residue class lifts uniquely to a root across the nilpotent ideal since the derivative is $-1$, and the roots are closed under addition and multiplication in characteristic two. We identify $k$ with this subfield. Define
\[
 S=\{v^2:v\in\m\}\subseteq W=\m^2,\qquad
 s=|S|,\qquad w(A)=\frac{s}{Q^2}.
\]
Squaring induces a Frobenius-semilinear map $\m/\m^2\to W$. Since $k$ is perfect, $S$ is a $k$-subspace. Thus $s=q^{\dim_k S}$, including $s=1$ in rank zero. Multiplication on $W$ by any element of $A$ factors through its residue. In particular $S$ is an ideal.

\begin{theorem}\label{th:main}
For every $L\ge2$ and every such $A$,
\[
        (1-1/q)^2w(A)\ \le\ 1-\Delta_L(A)\ \le\ w(A).
\]
For any sequence $A_i$ of these rings and $R_t=\prod_{i=1}^t A_i$,
\[
 \lim_{t\to\infty}\Delta_L(R_t)>0
       \quad\Longleftrightarrow\quad
              \sum_{i\ge1}w(A_i)<\infty .
\]
\end{theorem}

The one-factor estimates determine the first-order loss uniformly when the residue field grows. For fixed $q$ their constants remain independent of the dimensions of both nilpotent layers. The product criterion is exact even at fixed $q$.

\section{A square fixes the top two coefficients}

Let $\bar f$ denote reduction to $k[x]$. If $f$ and $r^2f$ have nonzero reductions and belong to $B_n(A)$, and $e=\deg\bar r\ge1$, then $d=\deg\bar f\le n-2e$. Write $r=t+v$, where $t=\bar r\in k[x]$ and $v\in\m[x]$. Characteristic two implies
\[
          r^2f=t^2f+v^2\bar f.
\]
The square $v^2$ has coefficients in $S$ and only even powers of $x$. Work at infinity with formal Laurent series, writing $\deg$ for the largest exponent of $x$. The inverse of $t^2$ has only even powers. Consequently $U=v^2/t^2$ has coefficients in $S$ and only even powers, and
\[
            \deg(f+U\bar f)\le n-2e.             \tag{1}
\]
Here division by a scalar polynomial is legitimate in $A((x^{-1}))$. Equation (1) gives more than the assertion that the top coefficients belong to $W$.

\begin{lemma}\label{lem:tops}
Fix the coefficients of $f$ up through degree $n-2$. If $f$ can be a first term of a pair $f,r^2f$ as above, the ordered pair $(f_n,f_{n-1})$ has at most $s$ possible values. More precisely, write $F=\bar f$, $d=\deg F$, and take $F_{-1}=0$. If $n-d$ is even the possibilities lie in
\[
              \{(cF_d,cF_{d-1}):c\in S\}.
\]
If $n-d$ is odd they lie in $\{(0,cF_d):c\in S\}$.
\end{lemma}
\begin{proof}
The residue degree bound shows that the already fixed lower coefficients determine $F$. It is nonzero. If $U\ne0$, its leading coefficient times $F_d$ is nonzero and hence $\deg(UF)=\deg U+d\le n$ by (1). Put $j=2\lfloor(n-d)/2\rfloor$. The top coefficients at $n,n-1$ of $UF$ can involve only $U_j$: every lower available index is at most $j-2$. In characteristic two they equal those of $f$ by (1). Taking $c=U_j$ proves the assertion, also when $U=0$.
\end{proof}

Delete from $B_n(A)$ the pairs of highest coefficients in Lemma~\ref{lem:tops} for each nonzero lower residue polynomial, with both upper residues zero as the formulas imply. Also delete all zero-reduction polynomials. The deleted fraction is at most $s/Q^2+q^{-(n+1)}$. In a remaining forbidden triple the reduction of the ratio must have degree $e\ge1$. Indeed zero reduction would kill a residue, and a nonzero constant reduction would make the ratio a unit by nilpotence. At least one of $a,b,a+b$ is even. The corresponding two terms therefore have the form $g,u^2g$ with $\deg\bar u\ge1$. The first of them was deleted.

\section{A matching of the same order}
We construct vertex-disjoint triples in $B_n(A)$, for $n\ge3$. Allow arbitrary coefficients $f_0,\ldots,f_{n-2}$ with $f_0$ and $f_{n-2}$ units. Let $F_j=\bar f_j$ and choose any $c\in S$. Complete the start $f$ by
\[
 f_n=cF_{n-2},\qquad f_{n-1}=cF_{n-3}.
\]
For every $c$, choose once and for all $v=v(c)\in\m$ with $v^2=c$, taking $v(0)=0$. Put
\[
 r=x+v(x^2+1)+cx^3,\qquad y=rf,\qquad z=r^2f .
\]
In $y$, the coefficient at $n+1$ is $f_n+cf_{n-2}=0$, since $vf_{n-1}=0$. Higher coefficients vanish because $\m W=0$. In $z$ use $r^2=x^2+c(x^4+1)$: the coefficients at $n+2,n+1$ are $f_n+cF_{n-2}$ and $f_{n-1}+cF_{n-3}$, and both vanish. All three terms belong to $B_n(A)$ and their residue degrees are $n-2,n-1,n$.

From either $y$ or $z$ we recover $\bar f$ by division by $x$ or $x^2$ in $k[x]$. For the end, $z(0)=cF_0$ recovers $c$. For the middle, $y(0)=v f_0$ determines $v\bmod W$ because $F_0\ne0$ and $\m/W$ is a $k$-vector space. This determines $v^2=c$. We now know the chosen lift $v(c)$ and $r$. Multiplication by $r$ and by $r^2$ is injective because $r^4=x^4$. Thus the triples are disjoint, with starts occupying the fraction
\[
                    u=(1-1/q)^2\frac{s}{Q^2}
\]
of the ball. Every admissible set misses at least one term per triple. Together with the deletion argument and the gluing step below, this proves the local bounds.

\section{Products and the infinite sets behind the bounds}
For a fixed product, the supremal upper density equals the upper limit of the maxima in the finite balls. Here are the details needed to pass from one to the other. Take cutoffs $n_j$ realizing that upper limit, successively large enough that $n_j-j>(1+L)n_{j-1}$ and $n_j-(1+L)j>n_{j-1}$. From each finite optimizer retain only points whose residue degree in every coordinate is between $n_j-j$ and $n_j$. This discards at most $\sum_i q_i^{-j-1}$ of the ball. Take the union of the retained blocks. Along a potential triple all residue degrees are finite and have the form $d,d+ae,d+(a+b)e$ in each coordinate, $e\ge0$. The block indices are nondecreasing. If only the last term belongs to the latest block, the last degree is at most $(1+L)n_{j-1}$. If the last two are in the latest block then $be\le j$, so the first degree is at least $n_j-(1+L)j$. Both exclude a triple crossing blocks. A triple within one block is already forbidden. This proves the equality.

Taking products of the survivors from the deletion argument gives
\[
                  \Delta_L(R_t)\ge\prod_{i\le t}(1-w(A_i)).
\]
Indeed, with all residues nonzero, a nonunit ratio in the product has positive residue degree in some coordinate. For the reverse bound, partition each coordinate ball into the disjoint triples constructed above and leftover singletons. For a uniformly chosen point the number $K$ of triple cells it meets is a sum of independent Bernoulli variables of parameters $3u_i$. Given its cells the point is uniform on a ternary cube of dimension $K$. A combinatorial line, moving through the three positions simultaneously in a nonempty set of coordinates, is a forbidden triple: use the specified ratio in each moving coordinate and $1$ elsewhere. If $\alpha_j$ denotes the largest density of a line-free set in $\{0,1,2\}^j$, $\alpha_0=1$, every finite maximum is at most $\mathbb E[\alpha_K]$.

Slicing with one coordinate held fixed shows that $\Delta_L(R_t)$ decreases with $t$. We have $u_i\ge w(A_i)/4$, $0<w(A_i)\le1/Q_i\le1/2$, and $\operatorname{Var}K\le\mathbb EK$. A convergent sum of the $w(A_i)$ gives a positive infinite product. For a divergent sum, $K$ tends to infinity in probability and $\alpha_j\to0$ by the density Hales--Jewett theorem \cite{dhj}. Thus $\mathbb E[\alpha_K]\to0$, proving Theorem~\ref{th:main}.

\section{Two multiplication laws with opposite answers}
Fix a field $k=\F_q$ of characteristic two and $h\ge1$. On $A=k\oplus V\oplus W$ take bases $e_0,\ldots,e_h$ of $V$ and $w_1,\ldots,w_h$ of $W$, let $1$ act as the identity, and set $(V\oplus W)W=0$. Specify the symmetric product on $V$ by $e_0e_j=w_j$ for $1\le j\le h$. All unspecified products are zero. Call this algebra $A_h^0$. Define $A_h^1$ by additionally setting $e_j^2=w_j$ for $1\le j\le h$. Both laws are associative since all products of three radical elements vanish. Both have $\m^2=W$ and the same radical-layer dimensions $(1,h+1,h)$. Nevertheless
\[
 \begin{array}{c|c|c|c}
  &Q&s&w(A)\\\hline
 A_h^0&q^{2h+2}&1&q^{-4h-4}\\
 A_h^1&q^{2h+2}&q^h&q^{-3h-4}
 \end{array}
\]
because the mixed terms in a square cancel. Set $h_i=\lceil(2/7)\log_q(i+1)\rceil$. Up to fixed positive factors the losses in the criterion are $(i+1)^{-8/7}$ and $(i+1)^{-6/7}$. The first sequence of products has positive limiting extremal density and the second has limit zero. The dimensions of the layers thus fail to determine even positivity of the limit.

\paragraph{Scope and remaining question.} The square-zero characteristic-two result in the repository reappears at the level of the product criterion when $s=1$. Its sharper constant for that special case comes from the direct $x$-shift packing. Here the new factor $s$ can grow without bound. The theorem does not determine each individual value of $\Delta_2(A)$, and the proof uses $2=0$ in the coefficient ring itself. Residue characteristic two with $2\ne0$ produces the additional term $2tv$ in $(t+v)^2$. An exact invariant there, or for $\m^4=0$, remains an open extension.

\begin{thebibliography}{9}
\bibitem{asada} M. Asada, E. Fourakis, S. Manski, N. McNew, S. J. Miller and G. Moreland,
\emph{Subsets of $\mathbb F_q[x]$ free of 3-term geometric progressions}, Finite Fields Appl. 44 (2017), 135--147. \href{https://arxiv.org/abs/1512.01932}{Preprint and Question 6.5}.
\bibitem{repo} ProveIt contributors, \emph{Beyond the Riemann hypothesis for Dirichlet $L$-functions}, section on progression-free sets in polynomial and finite rings, labels \texttt{gpp:thm:local} and \texttt{gpp:q:cube}. \href{https://github.com/VladimirReshetnikov/ProveIt/blob/ff96fd44ab160e65fb8a9d861243828d07c156cf/Analysis/GeneralizedRiemannHypothesis/Article/sections/gpp.tex}{Pinned article source}.
\bibitem{dhj} D. H. J. Polymath, \emph{A new proof of the density Hales--Jewett theorem}, Ann. of Math. 175 (2012), 1283--1327. \href{https://doi.org/10.4007/annals.2012.175.3.6}{doi:10.4007/annals.2012.175.3.6}.
\end{thebibliography}
\end{document}
